\documentclass[11pt,a4paper]{article}
\usepackage[utf8]{inputenc}
\usepackage[T1]{fontenc}
\usepackage{amsmath,amssymb,amsthm,physics,geometry,booktabs,array,hyperref,mathtools}
\geometry{margin=2.5cm}
\hypersetup{colorlinks=true,linkcolor=blue,citecolor=blue,urlcolor=blue}
\theoremstyle{plain}
\newtheorem{theorem}{Theorem}[section]
\theoremstyle{definition}
\newtheorem{definition}{Definition}[section]
\theoremstyle{remark}
\newtheorem{remark}{Remark}[section]
\newcommand{\lQ}{\ell_Q}

\title{\textbf{Quantum Gravitational Dynamics}\\[4pt]
\Large Chapter 4: Gravitational Energy}
\author{Romeo Matshaba}
\date{University of South Africa}

\begin{document}
\maketitle

\begin{abstract}
General Relativity possesses no well-defined local energy-momentum tensor for the gravitational field.  QGD resolves this century-old problem: the graviton field $\sigma_\mu$ admits a true Noether tensor $T^{\mu\nu}_{\mathrm{QGD}} = \nabla^\mu\sigma_\alpha\nabla^\nu\sigma^\alpha - \tfrac{1}{2}g^{\mu\nu}(\nabla\sigma)^2 + \lQ^2$-corrections, which is coordinate-independent, positive-definite, and locally conserved.  We derive the Schwarzschild energy profile, two-body binding energy, exact gravitational wave energy, and the gravitational Poynting vector.
\end{abstract}

\tableofcontents

%=============================================================
\section{The Problem in General Relativity}
\label{sec:problem}
%=============================================================

The gravitational field in GR carries energy --- gravitational waves heat up detectors --- yet the theory provides no local, covariant energy-momentum tensor for this energy.  The standard remedies are:

\begin{itemize}
  \item The Einstein pseudotensor $t_{\mu\nu}$: coordinate-dependent, can be set to zero at any point by choosing free-fall coordinates.
  \item The Landau--Lifshitz pseudotensor: symmetric but still non-covariant.
  \item The Bondi mass: defined only at null infinity; cannot locate energy in space.
  \item The ADM mass: defined only at spatial infinity.
\end{itemize}

The root cause is the equivalence principle: in free fall, the metric is locally flat and its first derivatives vanish.  Since the pseudotensors are built from $\partial g$, they vanish too.

\textbf{QGD resolution:} the graviton field $\sigma_\mu$ generates the metric but is not itself a metric component.  It cannot be gauged away, just as the Yang--Mills field strength $F_{\mu\nu}$ persists even when the potential $A_\mu$ is gauged.  This structural shift is what makes localizable energy possible.

%=============================================================
\section{The QGD Energy-Momentum Tensor}
\label{sec:tensor}
%=============================================================

\subsection{Derivation from Noether's theorem}

The QGD action for the $\sigma$-field is
\begin{equation}
  S[\sigma] = \int\!\dd^4x\left[\frac{1}{2}\partial_\alpha\sigma_\mu\partial^\alpha\sigma^\mu - V(\sigma) + \frac{\lQ^2}{2}(\partial_\alpha\partial_\beta\sigma_\mu)^2 + \mathcal{L}_{\mathrm{matter}}[g[\sigma]]\right]
\end{equation}

Applying Noether's theorem to translation invariance $x^\mu \to x^\mu + a^\mu$:
\begin{equation}
  \boxed{T^{\mu\nu}_{\mathrm{QGD}} = \nabla^\mu\sigma_\alpha\nabla^\nu\sigma^\alpha - \frac{1}{2}g^{\mu\nu}(\nabla\sigma)^2 - g^{\mu\nu}V(\sigma) + \lQ^2\text{-corrections}}
  \label{eq:Tmunu}
\end{equation}

\subsection{Properties}

\begin{enumerate}
  \item \textbf{True tensor:} transforms covariantly under all coordinate changes.  Cannot be set to zero by any gauge choice.
  \item \textbf{Conservation:} $\partial_\mu T^{\mu\nu}_{\mathrm{QGD}} = 0$ follows automatically from the field equations.  No pseudotensor gymnastics required.
  \item \textbf{Positive-definite Hamiltonian:}
  \begin{equation}
    H[\sigma,\pi] = \int\!\dd^3x\left[\frac{1}{2}\pi_\mu\pi^\mu + \frac{1}{2}(\nabla\sigma_\mu)^2 + V(\sigma) + \frac{\lQ^2}{2}(\nabla^2\sigma_\mu)^2\right] \geq 0
    \label{eq:H_positive}
  \end{equation}
  Equality holds if and only if $\sigma_\mu = 0$ (flat Minkowski).  This is manifest; it supersedes the 64-year effort culminating in the Schoen--Yau proof (1979).
\end{enumerate}

%=============================================================
\section{Local Energy Density}
\label{sec:density}
%=============================================================

The $T^{00}$ component gives the gravitational energy density at every spacetime point:
\begin{equation}
  \boxed{\rho_{\mathrm{grav}}(\mathbf{x}) = \underbrace{\frac{1}{2}\dot\sigma_\mu\dot\sigma^\mu}_{\text{kinetic}} + \underbrace{\frac{1}{2}(\nabla\sigma_\mu)^2}_{\text{gradient}} + \underbrace{V(\sigma)}_{\text{potential}} + \underbrace{\frac{\lQ^2}{2}(\nabla^2\sigma_\mu)^2}_{\text{quantum stiffness}}}
  \label{eq:rho}
\end{equation}

Each term is manifestly non-negative (for $V \geq 0$).  This is a well-defined scalar field at every point in spacetime.

%=============================================================
\section{Energy Density of the Schwarzschild Field}
\label{sec:schwarzschild}
%=============================================================

For a Schwarzschild black hole of mass $M$:
\begin{equation}
  \sigma_t(r) = \sqrt{\frac{r_s}{r}}, \qquad r_s = \frac{2GM}{c^2}
\end{equation}

The gradient energy density is:
\begin{equation}
  \rho_{\mathrm{grav}}(r) = \frac{1}{2}\!\left(\frac{\dd\sigma_t}{\dd r}\right)^{\!2} = \frac{1}{2}\cdot\frac{r_s}{4r^3} = \boxed{\frac{GM}{4c^2 r^3}}
  \label{eq:rho_schw}
\end{equation}

Physical interpretation:
\begin{itemize}
  \item Energy concentrated near the horizon: $\rho \propto r^{-3}$.
  \item At $r = r_s$: $\rho = c^4/(32G^2M^2)$.
  \item At $r = 10\,r_s$: three orders of magnitude smaller.
  \item No energy at the singularity: the quantum stiffness $\lQ^2(\nabla^2\sigma)^2$ activates at $r \sim \lQ$ and prevents unlimited concentration.
\end{itemize}

\subsection{Newtonian limit}

In the weak-field regime with Newtonian potential $\Phi = -GM/r$:
\begin{equation}
  \rho_{\mathrm{grav}} = \frac{1}{2}(\nabla\sigma_t)^2 = \frac{(\nabla\Phi)^2}{2c^2} \propto \frac{g^2}{8\pi G}
\end{equation}

This is exactly the classical gravitational field energy density, derived here from the $\sigma$-field structure.

%=============================================================
\section{Two-Body Binding Energy}
\label{sec:binding}
%=============================================================

For two masses $M_1$, $M_2$ separated by distance $d$, the $\sigma$-field superposition is:
\begin{equation}
  \sigma_t(\mathbf{x}) = \sqrt{\frac{2GM_1}{c^2|\mathbf{x}-\mathbf{x}_1|}} + \sqrt{\frac{2GM_2}{c^2|\mathbf{x}-\mathbf{x}_2|}}
\end{equation}

The total field energy splits as $E = E_1 + E_2 + E_{\mathrm{cross}}$, where the cross term from $\int\nabla\sigma_1\cdot\nabla\sigma_2\,\dd^3x$ gives:
\begin{equation}
  \boxed{E_{\mathrm{binding}} = -\frac{GM_1M_2}{d}}
  \label{eq:binding}
\end{equation}

Newton's gravitational binding energy emerges automatically from the field-gradient structure.

%=============================================================
\section{Gravitational Wave Energy}
\label{sec:gw}
%=============================================================

\subsection{Exact, point-wise definition}

A gravitational wave is a radiation mode $\sigma_\mu^{\mathrm{GW}} = \epsilon_\mu e^{ik\cdot x}$ with $k^2 = 0$.  The energy density is:
\begin{equation}
  \boxed{\rho_{\mathrm{GW}} = \frac{1}{2}(\partial\sigma^{\mathrm{GW}})^2 = \frac{1}{2}\omega^2|\epsilon|^2}
  \label{eq:rho_gw}
\end{equation}

This is exact, point-wise defined, and gauge-invariant.  No averaging over wavelengths required.

\subsection{Recovery of Isaacson}

In the weak-field limit, the LIGO observable gives $h_{ij}^{\mathrm{TT}} \sim -2\sigma_i\sigma_j$, so:
\begin{equation}
  \rho_{\mathrm{GW}} \;\longrightarrow\; \frac{c^2}{32\pi G}\langle\dot h^2\rangle
\end{equation}
recovering the Isaacson (1968) formula, which in GR requires an averaging procedure that is approximate and valid only in the short-wave limit.

\subsection{The gravitational Poynting vector}

Energy flux at any radius (not just infinity):
\begin{equation}
  \boxed{\mathbf{S}_{\mathrm{grav}} = \dot\sigma_\mu\nabla\sigma^\mu, \qquad \frac{\dd E}{\dd t} = \oint_S\mathbf{S}_{\mathrm{grav}}\cdot\dd\mathbf{A}}
  \label{eq:poynting}
\end{equation}

This enables energy-flow tracking from the near-zone orbital region through the wave zone to infinity --- a capability absent in GR.

%=============================================================
\section{The $Q_\mu$ Self-Energy Term}
\label{sec:Q}
%=============================================================

The master field equation $\Box_g\sigma_\mu = Q_\mu + G_\mu + T_\mu + \kappa\lQ^2\Box_g^2\sigma_\mu$ contains the self-interaction source:
\begin{equation}
  Q_\mu = \sigma_\mu(\nabla_\alpha\sigma_\beta\nabla^\alpha\sigma^\beta) + (\nabla_\mu\sigma_\alpha)(\sigma_\beta\nabla^\beta\sigma^\alpha)
\end{equation}

$Q_\mu$ is the gravitational field's self-energy current: the contribution from the field's own energy density.  The pattern across gauge theories:

\begin{table}[h]
  \centering
  \begin{tabular}{@{}lll@{}}
    \toprule
    Theory & Self-coupling & Consequence \\
    \midrule
    Maxwell ($Q = 0$) & None & Linear, no self-energy \\
    Yang--Mills ($Q \sim A\cdot F$) & Gluon self-interaction & Confinement \\
    QGD ($Q \sim \sigma(\partial\sigma)^2$) & Graviton self-energy & Black holes, GW energy \\
    \bottomrule
  \end{tabular}
\end{table}

%=============================================================
\section{Comparison: GR vs.\ QGD Energy}
\label{sec:comparison}
%=============================================================

\begin{table}[h]
  \centering
  \renewcommand{\arraystretch}{1.3}
  \begin{tabular}{@{}p{3.8cm}p{4.8cm}p{5cm}@{}}
    \toprule
    Property & GR & QGD \\
    \midrule
    Local energy density & Pseudotensor only & Scalar $\rho = \tfrac{1}{2}(\partial\sigma)^2$ \\
    True tensor & Does not exist & Noether $T^{\mu\nu}_{\mathrm{QGD}}$ \\
    Conservation & Coordinate-dependent & Automatic \\
    Positive energy & Schoen--Yau (1979, at $\infty$) & Manifest: $H \geq 0$ locally \\
    GW energy & Isaacson averaging & Exact: $\rho_{\mathrm{GW}} = \tfrac{1}{2}\omega^2|\epsilon|^2$ \\
    Black hole energy & ADM mass (at $\infty$) & Distributed: $\rho \propto r^{-3}$ \\
    Binding energy & Linearised GR; source undefined & $-GM_1M_2/r$ from gradients \\
    Energy flux & Only at null infinity & Poynting vector at any $r$ \\
    \bottomrule
  \end{tabular}
\end{table}

%=============================================================
\section{Conclusions}
\label{sec:conclusions}
%=============================================================

The century-long mystery of gravitational energy was a consequence of using the metric as the fundamental variable.  In QGD:

\begin{enumerate}
  \item A true energy-momentum tensor $T^{\mu\nu}_{\mathrm{QGD}}$ exists, derived by Noether's theorem.
  \item Energy density $\rho_{\mathrm{grav}} = \tfrac{1}{2}(\partial\sigma)^2$ is a well-defined scalar field at every point.
  \item Conservation is automatic; positive energy is manifest.
  \item Gravitational wave energy is exact and point-wise.
  \item Two-body binding energy $-GM_1M_2/r$ emerges from field-gradient cross terms.
  \item The gravitational Poynting vector tracks energy flow at any radius.
\end{enumerate}

\end{document}
